\chapter{Computational Systems}\label{ch:systems}

\section{Turing Machines}

A Turing machine $\mathrm{TM}(n,m)$ has $n$ states and $m$ colors (tape symbols).
The transition function maps (state, read symbol) to (new state, write symbol,
head direction).

\begin{definition}[Computational system]\label{comp-system}
\lean{CompSystem}
\leanok
A computational system $S$ consists of a configuration space $\mathrm{Config}_S$
and a partial step function $\mathrm{step}_S : \mathrm{Config}_S \to \mathrm{Config}_S$.
A configuration $c$ \emph{halts} if $\mathrm{step}_S^n(c) = \bot$ for some $n$.
\end{definition}

\begin{example}[Wolfram's $(2,3)$ TM]\label{wolfram-23-tm}
\notready
The smallest known universal Turing machine, with 2 states and 3 colors.
Smith~\cite{smith2007universality} proved its universality by simulating Rule~110.
\end{example}

\begin{example}[Rogozhin's small UTMs]\label{rogozhin-utms}
\notready
Rogozhin~\cite{rogozhin1996small} cataloged small universal TMs including $(7,4)$,
$(10,3)$, $(2,18)$, and others, establishing an efficiency frontier in the
$(n,m)$-plane.
\end{example}

\section{Cellular Automata}

An elementary cellular automaton (ECA) has 2 colors and nearest-neighbor rule on
$\ZZ$, parameterized by a rule number $r \in \{0,\ldots,255\}$.

\begin{example}[Rule 110]\label{rule-110}
\lean{ECA.rule110}
\leanok
Cook~\cite{cook2004universality} proved Rule~110 universal by simulating cyclic tag
systems via glider collisions. Neary and Woods~\cite{neary2006pcompleteness} further
showed Rule~110 is P-complete.
\end{example}

\begin{example}[Game of Life]\label{game-of-life}
\notready
Rendell~\cite{rendell2002turing} constructed an explicit Turing machine inside the
Game of Life, with spatial overhead $\sovh \approx 10^3$ cells and constant temporal
overhead.
\end{example}

\section{Tag Systems}

A \emph{tag system} $(m, \Sigma, P)$ has deletion number $m$, alphabet $\Sigma$, and
production rules $P : \Sigma \to \Sigma^*$. At each step: if the word is nonempty,
read the first symbol $a$, append $P(a)$, and delete the first $m$ symbols.

\begin{example}[Cyclic tag systems]\label{cyclic-tag}
\notready
A restricted variant where $\Sigma = \{0,1\}$, $m=1$, and the production rules cycle.
Cyclic tag systems are universal (they simulate 2-tag systems) and serve as an
intermediate in many simulation chains.
\end{example}

\section{Generalized Shifts}

C.~Moore's \emph{generalized shift}~\cite{moore1991generalized} is a stateless
machine on a bi-infinite tape. At each step a fixed-length window is read, rewritten,
and shifted. With window 3 and shift $\pm 1$, it coincides with Wolfram's
\emph{generalized mobile automaton} from NKS Chapter~11.

\section{Planar Billiards}

Miranda and Ramos~\cite{miranda2025billiards} proved that classical two-dimensional
billiard systems are Turing complete. For each TM $M$, they construct a planar
billiard table $B$ whose trajectories simulate $M$'s computation.

The encoding uses Topological Kleene Field Theory adapted to billiard flows:
\begin{enumerate}
\item The tape state $t$ and head position $k$ are encoded as a single point
  $x_{t,k} \in [0,1]$ via a Cantor-set-like ternary embedding: tape symbols
  become ternary digits, head position selects a nested subinterval.
\item The TM's finite state machine graph is transformed into a billiard table:
  each state $q_i$ corresponds to a region, each transition edge to a ``corridor.''
\item Shift (Lemma 1): parabolic walls implement $x_{t,k} \mapsto x_{t,k+1}$
  via affine maps on the interval.
\item Read-write (Lemma 2): walls with alternating slopes separate the 0/1
  subintervals and redirect trajectories to different corridors, implementing
  the transition function.
\end{enumerate}

Bennett's reversibility theorem is essential: the billiard construction requires
the TM to be reversible (injective transition function), and Bennett showed
every TM has a reversible equivalent.

\begin{remark}[Non-standard overhead]\label{billiard-overhead}
The TM $\to$ billiard edge has a non-standard overhead: the tape is encoded
as a single real number ($\sovh$ is not meaningful in the usual sense), and
$\tovh \approx \oh{|Q|}$ billiard reflections per TM step (one corridor
traversal with $\oh{|Q|}$ wall reflections). This edge demonstrates that
our $(\sovh, \tovh)$ framework needs extension for continuous dynamical
system vertices.
\end{remark}

\begin{remark}[Formalization difficulty]\label{billiard-formalization}
Lean formalization of this edge would be very hard: it requires Cantor set
analysis, billiard dynamics (specular reflection off piecewise-smooth boundaries),
parabolic mirror geometry, and Bennett's reversibility theorem. A realistic
partial target: formalize just the ternary encoding $t \mapsto x_t$ and the
shift map (Lemma 1). Long-term goal only.
\end{remark}

\section{Other Formalisms}

\begin{itemize}
\item \textbf{Register machines}: finite-state controller on unbounded integer
  registers. Simulate TMs with $\oh{1}$ overhead.
\item \textbf{SK combinators}: the combinatory logic basis $\{S,K\}$ simulates any
  TM with polynomial overhead.
\item \textbf{Lambda calculus}: Church-equivalent to TMs by standard results.
\end{itemize}
